\section{Recent constructions and the fixed-viscosity transfer problem}
\label{vh:source-scope}

\paragraph{Status and provenance (9 September 2026).}
This section is a dated research supplement, not a new premise of the
conditional regularity theorem. External author claims, the project's
acceptance of an input, independent review, and kernel verification are
separate evidence levels. The project source ledger is
\texttt{literature/recent-progress-2026-09.md}; its Euler inspection is
\texttt{literature/openai-euler-transfer-2026-09-09.md}.
Historical manuscript commits belong to this same project, not to an
independent prior presentation. No theorem in the canonical graph is
promoted by this supplement.

\paragraph{Forced Navier--Stokes is not the unforced target.}
OpenAI's September 8 manuscript states a bounded-energy singular solution
with smooth external force and zero initial velocity
\cite{OpenAINavier2026}. The project accepts that forced result as specified
in its live plan; the recorded kernel replication belongs to the pinned
source revision, not to a new build performed for this supplement.
It does not provide a singular trajectory with zero force and nonzero
Schwartz initial velocity. In particular a residual flat at the candidate
endpoint is not an identically vanishing solenoidal residual.

\paragraph{Unforced Euler and the preparation history.}
The separate Euler manuscript states a singular unforced solution from
smooth compactly supported solenoidal data \cite{OpenAIEuler2026}.
Its use here is as a proposed construction to adapt, not a theorem about
positive viscosity. The inspected ingredients include connected preparation
histories, a pressure-based displacement inverse, and summable initial
increments. At fixed positive viscosity one must change the actual
history equations and retain diffusion throughout preparation, rather than
switching it on only during the final growth interval. The publicly
served PDF gave different parsed and rendered pagination in this search;
we therefore do not assign a newly verified byte fingerprint to it.
The frozen source trail and the precise comparison are recorded in
\texttt{literature/viscous-history-source-audit-2026-09-09.md}.

\paragraph{Prior art for the pre-activation obstruction.}
Palasek's variable-gap Obukhov construction distinguishes an unforced
inviscid model from a viscous construction requiring a force
\cite{Palasek2026Obukhov}. Its discussion of the viscous case explicitly
tracks dissipation before a mode's activation, and uses the force to
remove that cost. This is prior art for the preparation problem, not an
original-NS embedding. We import neither the model's blowup nor its
unforced regularity into the PDE. The latter is attributed there to
work listed as ``to appear'', which was not independently inspected here.

\paragraph{Exact missing equation.}
For a globally defined solenoidal trial field $U$, put
\[
 \mathcal F_U=\partial_tU-\nu\Delta U+
                 \mathbb P\operatorname{div}(U\otimes U).
\]
An unforced correction must satisfy the entire equation
\[
 \partial_tw-\nu\Delta w+
 \mathbb P\operatorname{div}(U\otimes w+w\otimes U+w\otimes w)
       =-\mathcal F_U,
\]
with one compatible Schwartz initial trace and a proved
blowup-preserving error bound on the whole interval before the endpoint.
The final pressure is reconstructed from $U+w$; it need not equal a trial
pressure. No such correction, no endpoint-uniform inverse for it, and no
replacement arbitrary-data critical estimate is established here.

\subsection{A global upper-pressure-curvature obstruction to viscous transfer}
\label{pc:section}

\paragraph{Research status.}
The proofs in this subsection are author derivations, with independent
mathematical audit pending. They exclude a specified unforced viscous
transfer of the preceding Euler construction; they do not establish
arbitrary-data Navier--Stokes regularity. No claim of priority for a
pressure criterion or for the harmonic-analysis ingredients is made.
The source-inspection and test boundaries are recorded in
\texttt{research/evidence/2026-09-09-pressure-curvature-barrier.md}.

For a smooth classical branch, set
\[
 M(t)=\|u(t)\|_{L^\infty},\qquad
 K(t)=\max\{0,\sup_{x\in\R^3}\lambda_{\max}(D^2p(t,x))\},\qquad
 A(t)=\int_0^t\sqrt{K(s)}\,ds.
\]
Only the largest positive pressure eigenvalue is measured. In particular,
large negative eigenvalues are not bounded by hypothesis. Suprema over
rational spatial points and a countable dense set of unit directions give
measurability. On every compact classical interval these quantities are
finite. Any measurable nonnegative majorant for $K$ can be used instead.
The pressure is always the canonical whole-space double-Riesz pressure.

\begin{lemma}[One-sided curvature and mean oscillation]
\label{pc:semiconcavity}
Let $f\in C^2(\R^3)$, let
\[
 B=[f]_{\mathrm{BMO}}:=\sup_Q\frac{1}{|Q|}\int_Q|f-f_Q|<\infty,
\]
where $Q$ runs over Euclidean balls, and suppose $D^2f\le kI$ everywhere
for some $k\ge0$. Then
\[
 \|\nabla f\|_\infty\le48\sqrt{2kB}.
\]
\end{lemma}
\begin{proof}
Fix $x$ and $r>0$, write $Q=B_r(0)$, and set
$h(y)=k|y|^2/2-f(x+y)$. This is convex. Put
$b=\frac{1}{|Q|}\int_Q|h-h_Q|$. Convexity along segments from zero gives
$h(y/2)\le(h(0)+h(y))/2$. Averaging yields
\[
 h(0)\ge2h_{B_{r/2}}-h_Q\ge h_Q-16b.
\]
If $|z|\le r/2$, Jensen's inequality on $B_{r/2}(z)\subseteq Q$ gives
$h(z)\le h_Q+8b$. The supporting-plane inequality, with
$z=(r/2)\nabla h(0)/|\nabla h(0)|$ when the gradient is nonzero, therefore
implies $|\nabla h(0)|\le48b/r$. Furthermore,
$b\le B+kr^2/2$. Since $\nabla h(0)=-\nabla f(x)$,
\[
 |\nabla f(x)|\le48(B/r+kr/2).
\]
For $kB>0$ take $r=\sqrt{2B/k}$. If $k=0$ let $r\to\infty$;
if $B=0$ let $r\downarrow0$. These also handle the zero cases.
\end{proof}

\begin{lemma}[Canonical pressure oscillation]
\label{pc:bmo}
There is a fixed dimensional constant $C_p<\infty$ such that, for every
real solenoidal $u\in L^2\cap L^\infty$,
\[
 [R_iR_j(u_i u_j)]_{\mathrm{BMO}}\le C_p\|u\|_\infty^2.
\]
Here and below the repeated indices in the pressure are summed.
\end{lemma}
\begin{proof}
We give the needed endpoint singular-integral argument rather than assume
$L^\infty$ boundedness of the Riesz transforms, which is false. Write
$F=u_i u_j$ for one scalar tensor entry; $F\in L^1\cap L^\infty$ and
$\|F\|_\infty\le\|u\|_\infty^2$. For $Q=B_r(x_0)$ split
$F=F\mathbf1_{2Q}+F\mathbf1_{(2Q)^c}$. The Fourier multiplier
$-\xi_i\xi_j/|\xi|^2$ has absolute value at most one. Plancherel and
Cauchy--Schwarz consequently bound the mean absolute value of its near
output on $Q$ by $2^{3/2}\|F\|_\infty$.
Away from zero the kernel is $\partial_i\partial_j(4\pi|x|)^{-1}$;
its gradient is bounded by $C|x|^{-4}$. Any distribution supported at zero
in this kernel is a bounded local multiple of $F$ and is already covered
by the near part. For the far output subtract its value at $x_0$.
The mean-value estimate gives, for $x\in Q$,
\[
 |T(F\mathbf1_{(2Q)^c})(x)-T(F\mathbf1_{(2Q)^c})(x_0)|
 \le Cr\|F\|_\infty\int_{2r}^\infty \rho^{-4}\rho^2\,d\rho
 \le C'\|F\|_\infty.
\]
The subtracted value is well defined since $F\in L^1$. For any chosen
constant $c$, $\frac{1}{|Q|}\int_Q|TF-(TF)_Q|\le2\frac{1}{|Q|}\int_Q|TF-c|$.
Combining the estimates and summing the nine tensor entries proves the
claim. These estimates also justify the distributional splitting by
approximation. They are the classical endpoint argument; see
\cite{Stein1970} for its general setting.
\end{proof}

\begin{theorem}[Positive pressure-curvature clock]
\label{pc:clock}
There is a dimensional constant $C=48\sqrt{2C_p}$, independent of $\nu$,
such that every canonical unforced whole-space classical branch satisfies
\begin{equation}\label{pc:speed}
 M(t)\le M(0)e^{CA(t)}\qquad(0\le t<T_*).
\end{equation}
If $T_*<\infty$, necessarily
\begin{equation}\label{pc:necessary}
 \int_0^{T_*}e^{2CA(t)}\,dt=\infty.
\end{equation}
In particular $A(T_*)=\infty$. More generally, the finiteness of the
integral in \eqref{pc:necessary} implies continuation.
\end{theorem}
\begin{proof}
Lemmas~\ref{pc:semiconcavity}--\ref{pc:bmo} give at every classical time
\[
 \|\nabla p(t)\|_\infty\le C\sqrt{K(t)}M(t).
\]
For $v_\epsilon=(|u|^2+\epsilon^2)^{1/2}$, direct differentiation of the
original equation gives
\[
 (\partial_t+u\cdot\nabla-\nu\Delta)v_\epsilon
 =-\frac{u\cdot\nabla p}{v_\epsilon}
 -\nu\left(\frac{|\nabla u|^2}{v_\epsilon}
       -\frac{\sum_j(u\cdot\partial_j u)^2}{v_\epsilon^3}\right)
 \le |\nabla p|.
\]
The parenthesized term is nonnegative by Cauchy--Schwarz. The bounded
whole-space maximum principle on a compact classical time interval,
followed by $\epsilon\downarrow0$, implies
\[
 M(t)\le M(0)+\int_0^t\|\nabla p(s)\|_\infty\,ds.
\]
One may justify the whole-space principle by subtracting a vanishing
positive spatial barrier before taking its coefficient to zero. All
coefficients and sources are bounded on the compact interval. Integral
Gronwall proves \eqref{pc:speed}; no division by $M$ is needed.

For completeness the required viscous consumer can be proved directly.
Testing the equation with $-\Delta u$, using solenoidality to remove
pressure only in this pairing, and then Young's inequality gives
\[
 \frac{d}{dt}\|\nabla u\|_2^2+\nu\|\Delta u\|_2^2
 \le \nu^{-1} M(t)^2\|\nabla u\|_2^2.
\]
Thus $\int_0^{T_*}M(t)^2\,dt<\infty$ bounds the enstrophy and
contradicts the $H^1$ blowup alternative of
Proposition~\ref{prop:localtheory}. Equation~\eqref{pc:speed} proves
\eqref{pc:necessary}. If only $A(T_*)<\infty$ is used, there is also the
already reviewed endpoint chain
$M\hbox{ bounded}+\|u\|_2^2\le\|d\|_2^2
\Rightarrow\|u\|_3^3\le M\|d\|_2^2
\Rightarrow$ Theorem~\ref{thm:continuation}.
\end{proof}

\begin{corollary}[A one-sided type-I threshold]
\label{pc:typeI}
A finite endpoint is impossible if, on a terminal interval,
$K(t)\le c(T_*-t)^{-2}$ with $2C\sqrt c<1$.
Consequently a finite singular endpoint requires
\[
 \limsup_{t\uparrow T_*}(T_*-t)^2 K(t)\ge\frac1{4C^2}.
\]
The constant is not claimed sharp.
\end{corollary}
\begin{proof}
On such an interval $A(t)\le A(s)+\sqrt c\log((T_*-s)/(T_*-t))$.
Hence $M(t)^2\le C_s(T_*-t)^{-2C\sqrt c}$ is integrable, and the
preceding enstrophy argument applies. A strict limsup below the displayed
threshold supplies a $c$ with the required strict inequality.
\end{proof}

\begin{proposition}[Uniform full-state no-transfer bound]
\label{pc:no-transfer}
Fix $\nu>0$ and $H<\infty$. Consider any family of exact original
unforced solutions $u_j$ defined on $[0,t_j]$, $t_j\le H$, with
canonical pressures and real solenoidal Schwartz initial data $d_j$.
Suppose
\[
 \sup_j\|d_j\|_{H^3}\le B_3,\qquad
 \sup_j\|d_j\|_\infty\le B_\infty,\qquad
 \sup_j\int_0^{t_j}\sqrt{K_j(t)}\,dt\le B.
\]
Then $\sup_j\|\nabla u_j(t_j)\|_\infty<\infty$. In particular,
no such family can simultaneously have summable initial increments in
every $H^m$, bounded $t_j$, a uniform global upper bound
$D^2p_j\le K_+I$, and $\|\nabla u_j(t_j)\|_\infty\to\infty$.
No hypothesis on packet supports, overlaps, tails, or polarizations is used.
\end{proposition}
\begin{proof}
Here is a uniform estimate, avoiding any exchange of endpoints in a
continuous-dependence argument. Write $\mathcal H_3$ for the norm obtained
by summing the squared $L^2$ norms of all ordered derivatives of orders
zero through three; it is equivalent to $H^3$. The product estimates
\[
 \|D^k(u\otimes u)\|_2\le c_3 M\|D^k u\|_2,
 \qquad 0\le k\le3,
\]
hold with dimension-dependent constants. The only nontrivial intermediate
products follow from
\[
 \|Du\|_4^2\le cM\|D^2u\|_2,\qquad
 \|Du\|_6^2\le cM\|D^2u\|_3,\qquad
 \|D^2u\|_3^2\le c\|Du\|_6\|D^3u\|_2.
\]
Each follows by integrating one derivative by parts in the corresponding
power of the derivative norm and applying Holder. Spatial cutoffs and
regularization of the Euclidean derivative-tensor magnitude justify the
zero sets and the infinite domain. The last two inequalities imply
$\|Du\|_6\|D^2u\|_3\le cM\|D^3u\|_2$; Leibniz then gives the product
estimates. Differentiating the equation, using its divergence form, and
pairing gives
\[
 \frac{d}{dt}\mathcal H_3(u)^2
 \le \frac{c_3}{\nu}M(t)^2\mathcal H_3(u)^2.
\]
Pressure vanishes only in these solenoidal pairings; it was retained in
\eqref{pc:speed}. Consequently
\[
 \|\nabla u_j(t_j)\|_\infty
 \le c_s B_3\exp\left(\frac{c_3}{2\nu}
                  B_\infty^2 H e^{2CB}\right),
\]
where $c_s$ includes norm equivalence and $H^3\hookrightarrow W^{1,\infty}$.
The latter embedding follows also by Fourier Cauchy--Schwarz, since
$\int |\xi|^2(1+|\xi|^2)^{-3}\,d\xi<\infty$ in dimension three.
For the final assertion, summable increments give $B_3,B_\infty<\infty$;
$B=H\sqrt{\max(K_+,0)}$ is an admissible common clock bound.
\end{proof}

\paragraph{What is excluded and what is not.}
The Euler source explicitly maintains one global upper pressure-Hessian
bound over its exact approximating sequence \cite{OpenAIEuler2026}.
Proposition~\ref{pc:no-transfer} rules out preserving this invariant in
\emph{any} fixed-positive-viscosity unforced realization of that sequence,
even with arbitrary autonomous overlap and the entire exterior retained.
Merely repairing the pre-activation heat loss cannot overcome this
obstruction. It does not refute the Euler result: the velocity estimate
also holds inviscidly, but its derivative continuation uses $\nu>0$.
It does not refute a forced construction either: both the speed equation
and the canonical pressure input would have to be changed for forcing.
Bounds confined to one trajectory or a proper core do not satisfy the
global semiconcavity hypothesis. A redesigned viscous construction with
unbounded positive-curvature clock and a different history inverse is
not excluded.

\paragraph{Terminal and originality boundaries.}
Pressure regularity criteria and endpoint $L^\infty$ to BMO estimates are
prior art. Related primary records include Seregin and Sverak's
\emph{Navier--Stokes equations with lower bounds on the pressure}
(2002, DOI 10.1007/s002050200199), and Chae and Constantin's
\emph{Remarks on type I blow up for the 3D Euler equations and the 2D
Boussinesq equations} (arXiv:2103.10672). The latter uses directional
signed Hessian/strain quantities for Euler, not the criterion proved
here. The former's precise full theorem was not imported: only its
institutional abstract and publication metadata were accessible in this
retrieval. Novelty and independent verification of the present combination
remain undetermined. We have not bounded $A$ or the exponential-history
integral from arbitrary initial data. Thus this is a complete
construction-class obstruction and a conditional continuation estimate,
not a positive producer for the original unrestricted problem.

\subsection{The full viscous displacement action is not the Euler action}
\label{va:section}

\paragraph{Consumer and scope.}
The pressure-curvature barrier has ruled out retaining the Euler construction's
uniform global upper-Hessian budget in a singular viscous realization.
A separate question remains: can the Euler displacement coercivity argument
at least be transferred locally to Navier--Stokes? The following exact identity
and full-flow counterexample answer negatively for the unrestricted displacement
form. They do not exclude a viscosity-specific parabolic inverse or a more
restricted admissible displacement space. All results below are author proofs;
independent audit and priority remain undetermined.

Write $G_{ij}=\partial_j u_i$, $H_p=D^2p$ and $D_t=\partial_t+u\cdot\nabla$
on an actual smooth unforced NS solution at fixed $\nu>0$. The solenoidal
linearized operator is
\[
 \mathcal L_u w=\mathbb P(D_t w+Gw-\nu\Delta w).
\]
For a smooth solenoidal displacement $\eta$ compactly supported in
$(0,T)\times\R^3$, put $w=D_t\eta-G\eta$. This $w$ is solenoidal:
the two gradient contractions in $\operatorname{div}(D_t\eta-G\eta)$ cancel.
Define
\[
 \mathfrak q_\nu(\eta)=-\int_0^T\!\int_{\R^3}
       \eta\cdot\mathcal L_u(D_t\eta-G\eta),\qquad
 \mathfrak q_E(\eta)=\int_0^T\!\int_{\R^3}
       (|D_t\eta|^2-\eta\cdot H_p\eta).
\]
The subscript $E$ identifies the Euler-shaped form, even when its coefficients
come from an actual NS background. It is not a substitution of an Euler
trajectory for that background.

\begin{proposition}[Complete viscous displacement identity]
\label{va:identity}
With the definitions above,
\begin{align}
 \mathcal L_u(D_t\eta-G\eta)
 =\mathbb P\bigg(&D_t^2\eta+H_p\eta-\nu\Delta D_t\eta
       +\nu G\Delta\eta
       +2\nu\sum_j(\partial_jG)\partial_j\eta\bigg),
 \label{va:operator}
\end{align}
and
\begin{align}
 \mathfrak q_\nu(\eta)=\mathfrak q_E(\eta)+\nu\int_0^T\!\int_{\R^3}
 \bigg(&\sum_j\partial_j\eta\cdot G\partial_j\eta
 -\sum_{j,k}G_{kj}\partial_j\eta\cdot\partial_k\eta
 -\sum_j\eta\cdot(\partial_jG)\partial_j\eta\bigg).
 \label{va:action}
\end{align}
In particular the viscous correction has no automatic nonnegative sign.
\end{proposition}
\begin{proof}
The gradient of the original equation gives
$D_tG+G^2=\nu\Delta G-H_p$. Expanding $D_tw+Gw$ first produces
$D_t^2\eta+(H_p-\nu\Delta G)\eta$. Expanding
$-\nu\Delta(D_t\eta-G\eta)$ cancels the term
$-\nu(\Delta G)\eta$ and gives \eqref{va:operator}.
No pressure term was set to zero before this calculation.

Since $\eta$ is solenoidal, the outer Leray projection disappears only
in its $L^2$ pairing. Integration by parts in $D_t$ uses
$\operatorname{div}u=0$ and the zero temporal boundary values. For the
viscous part integrate first in space. The identity
\[
 \partial_jD_t\eta=D_t\partial_j\eta
                      +\sum_kG_{kj}\partial_k\eta
\]
shows that $\int\eta\cdot\Delta D_t\eta$ contributes
$-\int\sum_{j,k}G_{kj}\partial_j\eta\cdot\partial_k\eta$;
its material time derivative integrates to zero. Integration of
$-\int\eta\cdot G\Delta\eta$ gives
$\int\sum_j[\partial_j\eta\cdot G\partial_j\eta+
\eta\cdot(\partial_jG)\partial_j\eta]$.
Combining the remaining factor $-2$ gives \eqref{va:action}.
Compact support of the test field justifies all boundary removals.
\end{proof}

For a rank-one leading derivative $\nabla\eta=a\otimes\xi$, the principal
viscous integrand in \eqref{va:action} is
\begin{equation}\label{va:symbol}
 |\xi|^2a\cdot G a-|a|^2\xi\cdot G\xi.
\end{equation}
Only the symmetric part of $G$ enters. For transverse unit vectors in two
strain eigendirections it is their eigenvalue difference, not a nonnegative
viscous energy. The following theorem realizes its negative sign around a
\emph{finite-energy actual NS flow}; an infinite-energy affine solution is
not used as the counterexample.

\begin{theorem}[Failure of Euler-shaped coercivity in arbitrarily small flow]
\label{va:negative}
Fix $\nu>0$ and $\delta>0$. There exists a nonzero real solenoidal
$d\in C_c^\infty(\R^3)^3$, with $\|d\|_{H^3}<\delta$, whose original
unforced solution is globally smooth, and a fixed $T>0$, with the following
properties. Its canonical pressure satisfies $H_p\le K I$ on $[0,T]$ for
a finite $K$ with $KT^2/\pi^2\le1/2$. There are real solenoidal test fields
$\eta_N\in C_c^\infty((0,T)\times\R^3)^3$ such that
\[
 \|D_t\eta_N\|_{L^2_{t,x}}=1,\qquad
 \mathfrak q_E(\eta_N)\ge\tfrac12,\qquad
 \mathfrak q_\nu(\eta_N)\longrightarrow-\infty.
\]
Thus neither a positive lower bound by the Euler-shaped action nor a finite
lower bound depending only on $\|\eta\|_2$ and $\|D_t\eta\|_2$ holds for
this unrestricted viscous displacement form on that one background.
\end{theorem}
\begin{proof}
Let $G_0=\operatorname{diag}(-1,1,0)$ and choose a real smooth compactly
supported cutoff $\theta$ equal to one on $B_2$. Set
\[
 d_\epsilon=\epsilon\,\nabla\times
             \left[-\frac{\theta(x)}3\,x\times G_0x\right].
\]
The identity $\nabla\times(x\times G_0x)=-3G_0x$ implies
$d_\epsilon=\epsilon G_0x$ on $B_2$. Its full global datum is solenoidal
and compactly supported, not affine. Choose $\epsilon>0$ sufficiently small
for the requested $H^3$ bound and the following elementary global bootstrap.
Testing NS with $-\Delta u$ and using Sobolev and interpolation gives
\[
 |\langle(u\cdot\nabla)u,\Delta u\rangle|
 \le c\|u\|_3\|\Delta u\|_2^2,\qquad
 \|u\|_3\le c\|u\|_2^{1/2}\|\nabla u\|_2^{1/2}.
\]
For small enough $\epsilon$, energy and a continuity bootstrap make
$c\|u\|_3<\nu/2$ and keep the enstrophy nonincreasing. The local $H^1$
blowup alternative supplies global smoothness. This argument uses the fixed
positive viscosity; its smallness threshold can depend on $\nu$.

Let $X(t,a)$ be this actual solution's flow map, $F=D_aX$. On a sufficiently
short interval all needed derivatives of $X,F,F^{-1},u,p$ are bounded,
$\det F=1$, and $X(0,a)=a$. On $a\in B_1$ define
\[
 v(t,a)=F(t,a)e_1,\qquad \xi(t,a)=F(t,a)^{-T}e_2,
\]
and let $G$ in the next expression be evaluated at $(t,X(t,a))$.
At $t=0$,
\[
 L(t,a):=|\xi|^2v\cdot Gv-|v|^2\xi\cdot G\xi=-2\epsilon.
\]
Uniform continuity permits choosing a fixed $T>0$, independent of $N$,
so that $L\le-\epsilon$ on $[0,T]\times B_1$. Shrinking $T$ also ensures
$KT^2/\pi^2\le1/2$, where $K$ is a common nonnegative upper bound for the
actual pressure Hessian on this interval. Such a bound exists by local
classical regularity, without prescribing the pressure.

Choose nonzero real $\phi\in C_c^\infty(B_1)$ and
$\chi\in C_c^\infty((0,T))$, and start with the exact solenoidal seed
\[
 z_N(a)=\nabla_a\times
           \left[e_3\phi(a)\frac{\sin(Na_2)}N\right].
\]
Define the unnormalised test field by its entire transported value,
\[
 \widetilde\eta_N(t,X(t,a))=\chi(t)F(t,a)z_N(a).
\]
Pushforward by a volume-preserving diffeomorphism preserves divergence:
for a scalar test $f$, changing variables gives
$\int Fz_N\cdot\nabla_xf(X)\,da
 =\int z_N\cdot\nabla_a(f\circ X)\,da=0$.
Thus these are legitimate compactly supported solenoidal spacetime tests.
They are not claimed to be independently evolving NS velocities.
Since $D_t(Fz_N)=G Fz_N$, their $L^2$ and material-derivative $L^2$ norms
are bounded independently of $N$. All derivatives of the actual background
on the fixed transported support are bounded.

Differentiating the complete seed, including its envelope corrections,
gives uniformly on that support
\[
 \nabla_x\widetilde\eta_N
 =-N\chi(t)\phi(a)\sin(Na_2)\,v(t,a)\otimes\xi(t,a)+O(1).
\]
The last term of \eqref{va:action}, which includes $\nabla G$, is $O(N)$
after integration; so are cross terms with the $O(1)$ derivative remainder.
The Euler-shaped action is $O(1)$. Volume preservation and the elementary
oscillatory limit for $\sin^2(Na_2)$ now give the exact limiting sign
\begin{equation}\label{va:limit}
 \lim_{N\to\infty}N^{-2}\mathfrak q_\nu(\widetilde\eta_N)
 =\frac\nu2\int_0^T\!\int_{B_1}\chi(t)^2\phi(a)^2 L(t,a)\,da\,dt<0.
\end{equation}
This limit follows by writing $\sin^2=(1-\cos(2Na_2))/2$ and integrating
by parts against the fixed smooth compact coefficient. No Fourier mode of
the background or test has been deleted.

For every compact test, apply the one-dimensional Dirichlet Poincare
inequality to $\eta(t,X(t,a))$ for each $a$, and use $\det F=1$:
\[
 \|\eta\|_2^2\le(T/\pi)^2\|D_t\eta\|_2^2,\qquad
 \mathfrak q_E(\eta)\ge\tfrac12\|D_t\eta\|_2^2.
\]
The unnormalised tests have a strictly positive limiting $L^2$ norm,
as follows from the corresponding $\cos^2$ average and the nonsingularity
of $F$. Their material-derivative norms are consequently bounded above
and away from zero. Divide by those norms. The Poincare inequality and
\eqref{va:limit} give all the conclusions.
\end{proof}

\begin{theorem}[Complete semiboundedness classification on finite-energy backgrounds]
\label{va:classification}
Let $u$ be an actual smooth unforced NS solution at fixed $\nu>0$ on an open
classical interval $I$, with the canonical pressure and the finite-energy
regularity of \textup{LOCAL}. The following are equivalent:
\begin{enumerate}
\item $u$ vanishes on $I\times\R^3$.
\item There is a finite $C_I\ge0$ such that every real solenoidal
$\eta\in C_c^\infty(I\times\R^3)^3$ satisfies
\[
 \mathfrak q_\nu(\eta)\ge
 -C_I\bigl(\|\eta\|_{L^2_{t,x}}^2+\|D_t\eta\|_{L^2_{t,x}}^2\bigr).
\]
\end{enumerate}
In fact, on every nonzero such background there is a fixed compactly
contained subinterval and a sequence supported in its transported spatial
patch for which
\[
 \|D_t\eta_N\|_2=1,\qquad
 \mathfrak q_E(\eta_N)\ge\tfrac12,\qquad
 \mathfrak q_\nu(\eta_N)\longrightarrow-\infty.
\]
The constant in the failed bound is allowed to depend on the entire fixed
background. Thus this is not merely a lack of uniformity over varying data.
\end{theorem}
\begin{proof}
If $u=0$, its canonical pressure is zero and \eqref{va:action} gives
$\mathfrak q_\nu=\|\partial_t\eta\|_2^2\ge0$.
For the converse, suppose $u(t_0)$ is nonzero at an interior classical time.
The symmetric strain $S=(G+G^T)/2$ has trace zero and obeys
\[
 \int_{\R^3}|S|^2=\frac12\int_{\R^3}|\nabla u|^2.
\]
Indeed expansion leaves the cross term
$\int\partial_j u_i\partial_i u_j=0$ by solenoidality and integration by
parts. This is justified by the global $H^1$ regularity. If $S$ vanished
everywhere, $\nabla u$ would vanish, and the only constant $L^2$ field is zero.
Consequently $S(t_0,x_0)$ has two distinct eigenvalues at some point $x_0$.
Choose orthonormal vectors $a_0,\xi_0$ in its minimum and maximum
eigendirections. The coefficient \eqref{va:symbol} there equals
$\lambda_{\min}(S)-\lambda_{\max}(S)<0$.

Translate $x_0$ to the origin and choose a positively oriented orthonormal
coordinate system with first two directions $a_0,\xi_0$. Initialize the
actual flow map at $t_0$, rather than prescribing a new background flow.
By continuity there is a spatial ball and a fixed forward subinterval
$J\Subset I$ on whose transported tube
\[
 |F^{-T}\xi_0|^2(Fa_0)\cdot G(Fa_0)
       -|Fa_0|^2(F^{-T}\xi_0)\cdot G(F^{-T}\xi_0)\le-\sigma
\]
for some $\sigma>0$. The actual pressure Hessian is bounded on a larger
compact classical time interval; shrink $J$ so that its bound $K$ satisfies
$K|J|^2/\pi^2\le1/2$.

The complete curl seed, pushforward, oscillatory calculation and material
Poincare estimate in the proof of Theorem~\ref{va:negative} now apply
verbatim in these rotated coordinates, replacing $\epsilon$ by $\sigma$.
The fields vanish near the endpoints of $J$ and are extended by zero within
$I$. The limit \eqref{va:limit} is strictly negative; after normalization
$\mathfrak q_\nu\to-\infty$ while
$\|\eta_N\|_2^2+\|D_t\eta_N\|_2^2\le1+(|J|/\pi)^2$.
This contradicts every finite $C_I$ and proves both the classification and
the stated stronger form of failure.
\end{proof}

\paragraph{What this theorem actually eliminates.}
The full linearized parabolic Cauchy problem has not become ill posed.
An indefinite nonsymmetric displacement form can have other useful inverses.
The excluded assertion is transfer of Euler's pressure-only displacement
coercivity to the full NS operator on all solenoidal displacements. Adding
only initial or terminal quadratic boundary penalties cannot repair this
counterexample: the test fields vanish near both temporal endpoints.
The negative term is present on a single globally regular small-data
background; it is not caused by a blowup hypothesis, a remote imposed force,
an infinite-energy affine approximation or a discarded exterior.
A proof on a smaller displacement space must establish that its constraints
exclude these tests and that the required NS corrections remain in that space.

\subsection{Testing the critical-curvature escape rather than stopping at a no-go}
\label{va:hardy}

One can abandon a uniform upper pressure bound without automatically losing
the Euler-shaped interior form. For a fixed terminal time $T_0$ suppose only
$H_p(t,x)\le c(T_0-t)^{-2}I$. On compactly supported displacements with zero
endpoints in $[0,T']$, $T'<T_0$, the classical Hardy calculation gives
\[
 \mathfrak q_E(\eta)\ge(1-4c)\|D_t\eta\|_2^2.
\]
To check uniformity in $T'$, pull back by the volume-preserving flow and
extend the test by zero to $T_0$. For each component write
$\zeta(t)=\sqrt{T_0-t}\,h(t)$. Direct expansion yields
\[
 |\zeta'|^2-\frac{|\zeta|^2}{4(T_0-t)^2}
 =(T_0-t)|h'|^2-\tfrac12(|h|^2)'.
\]
Integration proves Hardy's constant four and the assertion. Hence $c<1/4$
is an available \emph{interior variational} range with a logarithmically
divergent majorant for the curvature clock. This is the classical Hardy
inequality, not a newly constructed pressure or NS solution. The preceding
NS type-I threshold already excludes sufficiently small $c$; no claim of
sharpness or of a nonempty family of singular flows in the remaining range
is made. Initial free-trace boundary conditions require their own estimates.

Most importantly, \eqref{va:action} remains the exact viscous form in this
range. The Hardy estimate does not control its spatial-derivative correction;
Theorem~\ref{va:negative} disproves such an automatic implication even when
$K$ is bounded and the Euler-shaped form is strongly coercive. The revised
common-trace construction therefore needs a genuinely viscous inverse, not
merely a different exponent in the Euler action. Testing the velocity
linearization with $w$ instead gives the exact identity
\[
 \tfrac12\partial_t\|w\|_2^2+\nu\|\nabla w\|_2^2
 =-\int w\cdot G w+\langle\mathcal L_u w,w\rangle.
\]
It restores ordinary dissipative control but leaves signed strain work.
No endpoint-uniform estimate for that work on a concentrating common history
has been produced here. Accordingly neither these identities nor the Hardy
repair supply the missing nonlinear inverse, a regenerative cascade,
a critical producer, or a resolution of NS-R3.
